% GENERATED FILE. Edit canonical lessons, then run npm run generate:books.
% canonical-source-hash: fnv1a64:5a59e159711395bf
% canonical-lessons: BN-C52-dhora, BN-C52-kata, BN-C52-thama, BN-C52-khoja, BN-C52-harano

\chapter{To hold, To cut, To stop, To look for, To lose}
\label{ch:bn-to-hold-to-cut-to-stop-to-look-for-to-lose}
\hlchaptermodality{52}

\begin{chapteropening}
\textbf{By the end of this chapter:} \emph{I can say to hold, to catch, to cut, to stop, to look for and to lose.}

Complete the last lesson of chapter 52: I can say to hold, to catch, to cut, to stop, to look for and to lose.
\end{chapteropening}

\section[\texorpdfstring{\bn{ধরা} (dhôrā)}{ধরা (dhôrā)}]{\bn{ধরা} (dhôrā) — to hold, to catch}

\label{lesson:BN-C52-dhora}



\begin{quote}
Before the new one: say the Bengali for to bring, then the Bengali for to count.
\end{quote}

\subsection*{You'll want to know: \bn{ধরা}}
\textbf{\bn{ধরা}} — \emph{dhôrā} — \textquotedblleft{}to hold, to catch\textquotedblright{}.

\begin{cousinweb}
\textbf{\bn{ধরা}} (\emph{dhôrā}) is a tadbhava, a word inherited from Sanskrit and reshaped by sound change. It comes from \emph{dharati}, \textquotedblleft{}holds, keeps\textquotedblright{} (CDIAL 6747).
\end{cousinweb}

\begin{grammarlens}[title={What you've built}]
One more word for this chapter.
\end{grammarlens}

\subsection*{Your turn}
\begin{itemize}
\raggedright
  \item \emph{Say it:} \emph{dhôrā}
  \item \emph{Say it:} \emph{dhôrā}, once more
  \item \emph{Say it:} say \emph{dhôrā}
  \item \emph{Recall:} say the Bengali for to bring, then the Bengali for to count, then say \emph{dhôrā} again
\end{itemize}

\begin{tcolorbox}[breakable,colback=teal!4,colframe=teal!35!black,title={Before you move on}]
What does \bn{ধরা} mean? (\textquotedblleft{}To hold, to catch\textquotedblright{}.) Say it once more.
\end{tcolorbox}
\section[\texorpdfstring{\bn{কাটা} (kāṭā)}{কাটা (kāṭā)}]{\bn{কাটা} (kāṭā) — to cut}

\label{lesson:BN-C52-kata}



\begin{quote}
Before the new one: say the Bengali for to count, then the Bengali for to hold.
\end{quote}

\subsection*{You'll want to know: \bn{কাটা}}
\textbf{\bn{কাটা}} — \emph{kāṭā} — \textquotedblleft{}to cut\textquotedblright{}.

\begin{cousinweb}
\textbf{\bn{কাটা}} (\emph{kāṭā}) is a tadbhava, a word inherited from Sanskrit and reshaped by sound change. It comes from \emph{kartati}, \textquotedblleft{}cuts\textquotedblright{} (CDIAL 2854).
\end{cousinweb}

\begin{grammarlens}[title={What you've built}]
One more word for this chapter.
\end{grammarlens}

\subsection*{Your turn}
\begin{itemize}
\raggedright
  \item \emph{Say it:} \emph{kāṭā}
  \item \emph{Say it:} \emph{kāṭā}, once more
  \item \emph{Say it:} say \emph{kāṭā}
  \item \emph{Recall:} say the Bengali for to count, then the Bengali for to hold, then say \emph{kāṭā} again
\end{itemize}

\begin{tcolorbox}[breakable,colback=teal!4,colframe=teal!35!black,title={Before you move on}]
What does \bn{কাটা} mean? (\textquotedblleft{}To cut\textquotedblright{}.) Say it once more.
\end{tcolorbox}
\section[\texorpdfstring{\bn{থামা} (thāmā)}{থামা (thāmā)}]{\bn{থামা} (thāmā) — to stop}

\label{lesson:BN-C52-thama}



\begin{quote}
Before the new one: say the Bengali for to hold, then the Bengali for to cut.
\end{quote}

\subsection*{You'll want to know: \bn{থামা}}
\textbf{\bn{থামা}} — \emph{thāmā} — \textquotedblleft{}to stop\textquotedblright{}.

\begin{cousinweb}
\textbf{\bn{থামা}} (\emph{thāmā}) is a tadbhava, a word inherited from Sanskrit and reshaped by sound change. It comes from \emph{stambhate}, \textquotedblleft{}supports, arrests\textquotedblright{} (CDIAL 13683).
\end{cousinweb}

\begin{grammarlens}[title={What you've built}]
One more word for this chapter.
\end{grammarlens}

\subsection*{Your turn}
\begin{itemize}
\raggedright
  \item \emph{Say it:} \emph{thāmā}
  \item \emph{Say it:} \emph{thāmā}, once more
  \item \emph{Say it:} say \emph{thāmā}
  \item \emph{Recall:} say the Bengali for to hold, then the Bengali for to cut, then say \emph{thāmā} again
\end{itemize}

\begin{tcolorbox}[breakable,colback=teal!4,colframe=teal!35!black,title={Before you move on}]
What does \bn{থামা} mean? (\textquotedblleft{}To stop\textquotedblright{}.) Say it once more.
\end{tcolorbox}
\section[\texorpdfstring{\bn{খোঁজা} (khõjā)}{খোঁজা (khõjā)}]{\bn{খোঁজা} (khõjā) — to look for}

\label{lesson:BN-C52-khoja}



\begin{quote}
Before the new one: say the Bengali for to cut, then the Bengali for to stop.
\end{quote}

\subsection*{You'll want to know: \bn{খোঁজা}}
\textbf{\bn{খোঁজা}} — \emph{khõjā} — \textquotedblleft{}to look for\textquotedblright{}.

\begin{cousinweb}
\textbf{\bn{খোঁজা}} (\emph{khõjā}) is a tadbhava, a word inherited and reshaped by sound change. It comes from \emph{*khojjati}, \textquotedblleft{}tracks, seeks\textquotedblright{}, a form CDIAL reconstructs from its descendants (CDIAL 3929). CDIAL doubts a link with Sanskrit \emph{khojati}, \textquotedblleft{}steals\textquotedblright{}.
\end{cousinweb}

\begin{grammarlens}[title={What you've built}]
One more word for this chapter.
\end{grammarlens}

\subsection*{Your turn}
\begin{itemize}
\raggedright
  \item \emph{Say it:} \emph{khõjā}
  \item \emph{Say it:} \emph{khõjā}, once more
  \item \emph{Say it:} say \emph{khõjā}
  \item \emph{Recall:} say the Bengali for to cut, then the Bengali for to stop, then say \emph{khõjā} again
\end{itemize}

\begin{tcolorbox}[breakable,colback=teal!4,colframe=teal!35!black,title={Before you move on}]
What does \bn{খোঁজা} mean? (\textquotedblleft{}To look for\textquotedblright{}.) Say it once more.
\end{tcolorbox}
\section[\texorpdfstring{\bn{হারানো} (hārāno)}{হারানো (hārāno)}]{\bn{হারানো} (hārāno) — to lose}

\label{lesson:BN-C52-harano}



\begin{quote}
Before the new one: say the Bengali for to stop, then the Bengali for to look for.
\end{quote}

\subsection*{You'll want to know: \bn{হারানো}}
\textbf{\bn{হারানো}} — \emph{hārāno} — \textquotedblleft{}to lose\textquotedblright{}.

\begin{cousinweb}
\textbf{\bn{হারানো}} (\emph{hārāno}) is a tadbhava, a word inherited from Sanskrit and reshaped by sound change. It comes from \emph{hārayati}, \textquotedblleft{}causes to be taken; loses\textquotedblright{} (CDIAL 14061). \emph{Hārāno} is the causing form of \emph{hārā}, \textquotedblleft{}to lose\textquotedblright{}, the verb CDIAL lists.
\end{cousinweb}

\begin{grammarlens}[title={What you've built}]
One more word for this chapter.
\end{grammarlens}

\subsection*{Your turn}
\begin{itemize}
\raggedright
  \item \emph{Say it:} \emph{hārāno}
  \item \emph{Say it:} \emph{hārāno}, once more
  \item \emph{Say it:} say \emph{hārāno}
  \item \emph{Recall:} say the Bengali for to stop, then the Bengali for to look for, then say \emph{hārāno} again
\end{itemize}

\begin{tcolorbox}[breakable,colback=teal!4,colframe=teal!35!black,title={Before you move on}]
What does \bn{হারানো} mean? (\textquotedblleft{}To lose\textquotedblright{}.) Say it once more.
\end{tcolorbox}
